% !TeX root = ../regression_logistique.tex
\chapter{Probabilité d'appartenance}
\label{ch:probabilite}

\orient{%
énoncer l'hypothèse de modélisation que pose la régression logistique ;
appliquer la fonction logistique à un score et interpréter le résultat ;
comparer sigmoïde, seuil, fonction affine par morceaux, tangente hyperbolique
et probit ; calculer une probabilité de bout en bout à partir de notes brutes}{%
exponentielle et logarithme (\annexeref{ann:prerequis}) ;
chapitre~\ref{ch:nuage}}{%
essentiel ; propriétés de $\sigma$ à l'\annexeref{ann:sigmoide}}

\section{Cote, logit et hypothèse du modèle}

Le score du chapitre précédent situe une observation par rapport à la frontière.
La régression logistique pose que ce score est le logarithme d'une cote.

La \emph{cote} d'un événement est le rapport entre sa probabilité et celle de
l'événement contraire :
\[
  \operatorname{cote}(p)
  =\frac{\annup{accent}{p}{\annl{probabilité}\\ \annl{de l'événement}}}
        {\ann{warning}{1-p}{\annl{probabilité}\\ \annl{du contraire}}}
\]
C'est la quantité employée dans les paris. Un élève auquel on accorde $p=0{,}8$
d'être \Gryf a ainsi une cote de
\[
  \frac{0{,}8}{0{,}2}=4 ,
\]
soit quatre fois plus de chances d'en être que de ne pas en être. Une cote de $1$
correspond à $p=0{,}5$, les deux réponses étant alors également probables. Quand
$p$ parcourt $]0,1[$, la cote parcourt $]0,+\infty[$ : elle tend vers $0$ lorsque
la probabilité s'annule, et croît sans borne lorsqu'elle approche $1$.

Le \emph{logit} est le logarithme de cette cote,
\[
  \operatorname{logit}(p)=\log\frac{p}{1-p}.
\]
Il vaut $0$ pour $p=0{,}5$, il est négatif en dessous et positif au-dessus, et il
parcourt $\R$ tout entier, comme le score. L'hypothèse de modélisation s'écrit
donc
\[
  \log\frac{p}{1-p}=z=w_0+w_1x_1+w_2x_2 .
\]
Le logit de la probabilité cherchée est supposé être une fonction \emph{affine}
des notes, c'est-à-dire une somme de leurs multiples augmentée d'une constante.
L'expression de $p$ en fonction de $z$ s'en déduit (\annexeref{ann:logit}).

Cette égalité est une hypothèse de modélisation. Que le logit et le score
parcourent tous deux $\R$ la rend \emph{possible} ; d'autres fonctions
conviendraient (ch.~\ref{ch:nuage}). Le choix du logit tient à trois
propriétés : il donne aux coefficients une lecture en log-cotes
(ch.~\ref{ch:modele}), il conduit à la vraisemblance de Bernoulli et à un coût
convexe (\annexeref{ann:cout}), et il produit un gradient simple
(\annexeref{ann:gradient}). C'est le lien canonique du modèle linéaire
généralisé binomial~\cite{mccullagh1989}.

Deux autres hypothèses interviennent ailleurs : l'indépendance conditionnelle des
observations, employée pour écrire la vraisemblance comme un produit
(\annexeref{ann:cout}), et le choix des variables retenues, qui fixe ce que le
modèle peut voir.

\begin{releve}
L'élève $3$, dont le score vaut $+1{,}34$, a une cote de $e^{1{,}34}=3{,}82$
d'être \Gryf, soit une probabilité de $3{,}82/4{,}82=0{,}79$.
\end{releve}

\section{Fonction logistique}

L'hypothèse se lit dans les deux sens : elle donne $z$ à partir de $p$, et c'est
$p$ à partir de $z$ dont le programme a besoin. L'inversion procède en quatre
opérations réversibles.

\begin{align*}
  \log\frac{p}{1-p} &= z \\[2pt]
  \frac{p}{1-p} &= e^{z}
    &&\text{exponentielle des deux membres} \\[2pt]
  p &= e^{z}(1-p) = e^{z}-p\,e^{z}
    &&\text{on multiplie par } 1-p \\[2pt]
  p\bigl(1+e^{z}\bigr) &= e^{z}
    &&\text{on regroupe les termes en } p \\[2pt]
  p &= \frac{e^{z}}{1+e^{z}} = \frac{1}{1+e^{-z}}
    &&\text{on divise en haut et en bas par } e^{z}
\end{align*}

La fonction ainsi obtenue s'appelle la \emph{fonction logistique}, et se note
$\sigma$ :
\[
  \ann{accent}{z}{\annl{score}\\ \annl{parcourt $\R$}}
  \;\longmapsto\;
  \ann{warning}{\sigma(z)=\frac{1}{1+e^{-z}}}{\annl{probabilité}\\ \annl{dans $\,]0,1[$}}
\]

Ses valeurs se déduisent de l'écriture $1/(1+e^{-z})$. Pour $z$ très négatif,
$e^{-z}$ est immense et le quotient tend vers zéro ; pour $z$ très positif,
$e^{-z}$ tend vers zéro et le quotient vers $1$ ; en $z=0$, l'exponentielle vaut
$1$ et le quotient exactement $0{,}5$. Sur la frontière, les deux réponses
reçoivent donc la même probabilité.

\begin{figure}[htbp]
\centering
\begin{graphique}{Courbe de la fonction logistique, en S, de 0 à 1, avec le score en abscisse de -6,5 à 6,5. Quatre points marqués : 0,5 en zéro, 0,79 pour un score de 1,34, 0,99 pour 4,71, et 0,02 pour -4,15.}
\begin{tikzpicture}
\begin{axis}[courseplot,height=6.4cm,axis lines=middle,
  xlabel={score $z$},ylabel={probabilité $\sigma(z)$},
  xmin=-6.5,xmax=6.5,ymin=-.12,ymax=1.15,ytick={0,.25,.5,.75,1},
  samples=200,domain=-6.5:6.5]
  \addplot[gridgray,thick,dashed,domain=-6.5:6.5]{1};
  \addplot[accent,very thick]{1/(1+exp(-x))};
  \addplot[only marks,mark=*,warning,mark size=2.6pt]
    coordinates {(0,0.5) (1.339,0.7922) (4.706,0.9910) (-4.154,0.0155)};
  \node[warning,anchor=west,font=\small] at (axis cs:1.6,.74){$0{,}79$};
  \node[warning,anchor=west,font=\small] at (axis cs:4.9,.93){$0{,}99$};
  \node[warning,anchor=east,font=\small] at (axis cs:-4.4,.12){$0{,}02$};
\end{axis}
\end{tikzpicture}
\end{graphique}
\caption{Fonction logistique. Les trois points marqués sont des élèves du
fichier, dont les scores figurent au tableau~\ref{tab:trois-eleves}.}
\label{fig:sigmoide-emploi}
\end{figure}

La fonction est croissante --- classer par score ou par probabilité revient au
même --- et symétrique autour de $0{,}5$, $\sigma(-z)$ valant $1-\sigma(z)$
(\annexeref{ann:sigmoide}).

\section{Fonctions de liaison et alternatives}

Toute transformation croissante centrée en $p=0{,}5$ conserve l'ordre des
scores et la frontière $z=0$. Le choix de la fonction porte donc sur le sens
probabiliste, la régularité de la perte et le comportement loin de la frontière,
tandis que la forme de la frontière provient du score.

\begin{table}[htbp]
\centering
\small
\begin{tabularx}{\textwidth}{
  >{\raggedright\arraybackslash}p{31mm}
  >{\raggedright\arraybackslash}p{40mm}
  >{\raggedright\arraybackslash}X}
\toprule
transformation & propriété & conséquence pour l'ajustement \\
\midrule
seuil $\mathbf{1}_{z\geq0}$ &
  sortie binaire, discontinuité en zéro &
  dérivée nulle ailleurs qu'au seuil : la sortie indique une décision, mais
  aucune direction de correction \\
affine écrêtée
  $\min(1,\max(0,\tfrac12+az))$ &
  probabilité par morceaux, angles et plateaux &
  gradient nul dans les zones écrêtées et rupture aux angles ; une prédiction
  assurée incorrecte peut cesser de fournir une correction exploitable \\
logistique $\sigma(z)$ &
  lisse, symétrique, valeurs dans $]0,1[$ &
  vraisemblance de Bernoulli, coût convexe pour un score affine et gradient
  $p-y$ \\
tangente hyperbolique &
  $\tfrac12[1+\tanh(z/2)]=\sigma(z)$ &
  même modèle après changement d'échelle ; les étiquettes $-1,+1$ donnent une
  notation équivalente \\
probit $\Phi(z)$ &
  fonction de répartition normale, lisse &
  autre modèle probabiliste valide, à queues différentes, ajusté lui aussi par
  vraisemblance \\
\bottomrule
\end{tabularx}
\caption{Transformations possibles d'un score réel. Le seuil sert à décider ;
les fonctions lisses servent à estimer et à ajuster une probabilité.}
\label{tab:liens}
\end{table}

Le lien complémentaire log-log fournit une autre option probabiliste, cette
fois asymétrique. Pour des classes mutuellement exclusives, la régression
\emph{softmax} remplace les sigmoïdes indépendantes par une distribution dont
les composantes somment à $1$. Une machine à vecteurs de support optimise une
marge avec une perte charnière ; elle produit d'abord un score de décision,
dont la conversion en probabilité demande une calibration séparée.

Une fonction de liaison différente modifie les probabilités et parfois
l'efficacité statistique, mais conserve une frontière affine tant que le score
reste affine. Des interactions, des fonctions polynomiales, un noyau, un arbre
ou un réseau peuvent produire une frontière courbe ; leur souplesse
supplémentaire augmente aussi les besoins en données, en réglage et en contrôle
du surapprentissage.

\section{Calcul de la probabilité}

La figure~\ref{fig:chaine-complete} associe la standardisation du
chapitre~\ref{ch:standardisation}, le score du chapitre~\ref{ch:nuage} et la
relation $p=\sigma(z)$. Ses constantes proviennent du fichier et de
l'entraînement.

\begin{chiffres}[Constantes du modèle]
\begin{center}
\begin{tabular}{lrrl}
\toprule
 & \Herb & \Runes & origine \\
\midrule
moyenne $\mu$      & $1{,}189$   & $495{,}052$  & le fichier \\
écart-type $\sigma$ & $5{,}175$   & $105{,}186$  & le fichier \\
\addlinespace
coefficient $w_j$  & $-3{,}454$  & $+3{,}469$   & l'entraînement \\
biais $w_0$        & \multicolumn{2}{c}{$-4{,}745$} & l'entraînement \\
\bottomrule
\end{tabular}
\end{center}
Les deux coefficients sont de signes opposés et de tailles voisines : une note
basse en \Herb et une note haute en \Runes poussent toutes deux le score vers
le haut.
\end{chiffres}

\begin{table}[htbp]
\centering
\small
\setlength{\tabcolsep}{5pt}
\begin{tabular}{llrrrrrr}
\toprule
élève & classe observée & \textcolor{herbo}{Herbo.} & $x_1$ & \textcolor{runes}{A.\ Runes} & $x_2$ & $z$ & $p$ \\
\midrule
$3$ & \Gryf & $-6{,}497$ & $-1{,}485$ & $523{,}982$ & $+0{,}275$ & $+1{,}34$ & $0{,}792$ \\
$4$ & \Gryf & $-7{,}82$ & $-1{,}741$ & $599{,}33$ & $+0{,}991$ & $+4{,}71$ & $0{,}991$ \\
$1$ & \Slyth  & $-5{,}99$ & $-1{,}387$ & $367{,}76$ & $-1{,}210$ & $-4{,}15$ & $0{,}016$ \\
$0$ & \Rav  & $+5{,}73$ & $+0{,}877$ & $532{,}48$ & $+0{,}356$ & $-6{,}54$ & $0{,}001$ \\
\bottomrule
\end{tabular}
\caption{Quatre élèves du fichier, de la note brute à la probabilité d'être
\Gryf.}
\label{tab:trois-eleves}
\end{table}

\begin{figure}[htbp]
\centering
\begin{graphique}{Trois barres horizontales à l'échelle, partant d'un axe vertical à zéro. Le biais s'étend vers la gauche jusqu'à -4,745 ; les deux termes de notes vers la droite, +5,129 et +0,954. Une quatrième barre, plus sombre, donne leur somme : +1,34.}
\begin{tikzpicture}[font=\small,x=0.85cm]
% Contributions signees au score de l'eleve 3, 1 unite = 1 cm.
\draw[ink!30] (0,-0.5) -- (0,3.6);
\node[font=\scriptsize,text=ink!70,anchor=south] at (0,3.62) {$0$};
\foreach \y/\lab/\val/\col in {%
  3.0/{$w_0\cdot 1=-4{,}745$}/-4.745/warning,
  2.0/{$w_1x_1=+5{,}129$}/5.129/accent,
  1.0/{$w_2x_2=+0{,}954$}/0.954/accent}{
  \fill[\col!65] (0,\y-0.22) rectangle (\val,\y+0.22);
  \node[font=\footnotesize,text=ink,anchor=west] at (5.6,\y) {\lab};
}
\draw[ink!30,dashed] (0,0.6) -- (0,0.2);
\fill[ink] (0,0.2) rectangle (1.338,-0.2);
\node[font=\footnotesize,text=ink,anchor=west] at (5.6,0) {$z=+1{,}34$};
\foreach \x in {-6,-4,-2,2,4}{
  \draw[ink!25] (\x,-0.45) -- (\x,-0.55);
  \node[font=\scriptsize,text=ink!70,anchor=north] at (\x,-0.55) {$\x$};
}
\draw[ink!40] (-6.4,-0.5) -- (5.4,-0.5);
\end{tikzpicture}
\end{graphique}
\caption{Les trois termes du score de l'élève $3$, à l'échelle : $+1{,}34$ pour
des contributions individuelles de l'ordre de $5$.}
\label{fig:contributions}
\end{figure}
\FloatBarrier

\begin{releve}
Ses notes standardisées valent $x_1=-1{,}485$ et $x_2=+0{,}275$, d'où
\[
  z=-4{,}745+(-3{,}454)\times(-1{,}485)+3{,}469\times0{,}275=+1{,}34,
  \qquad
  p=\frac{1}{1+e^{-1{,}34}}=0{,}792 .
\]
\end{releve}

Chaque terme du score est un produit de deux nombres signés, et le signe du
produit fixe le sens de la contribution : positif vers \Gryf, négatif vers les
autres classes. Le signe du coefficient est une propriété de la variable, fixée
par l'entraînement ; celui de la note standardisée situe l'observation par
rapport à la moyenne de sa colonne. Le biais entre sans facteur, contribution
identique pour toutes les observations, ici vers les autres classes : une
observation est classée \Gryf lorsque les deux termes de notes l'emportent sur
lui.

L'élève $3$ a une note d'\Herb basse et un coefficient négatif, donc un terme
positif ; sa note d'\Runes, légèrement au-dessus de la moyenne, contribue dans le
même sens. Les deux l'emportent de peu sur le biais : $z=+1{,}34$, $p=0{,}79$,
soit une position proche de la frontière. L'élève $4$ cumule une note d'\Herb
très basse et une note d'\Runes élevée, d'où $z=+4{,}71$ et $p=0{,}99$.

\begin{resultat}[Apport de la probabilité]
Deux observations rangées dans la même classe peuvent l'être à des distances très
différentes de la frontière. La probabilité chiffre cet écart là où la décision
binaire les confond, et c'est cette quantité continue qui permet, au chapitre
suivant, de mesurer de combien la frontière est mal placée.
\end{resultat}

\begin{exercice}[Probabilité de bout en bout]
L'élève $21$ a $+4{,}12$ en \Herb et sa note d'\Runes est absente. Standardiser
ses deux notes avec les constantes ci-dessus, calculer $z$ puis $p$, et dire de
quel côté de la frontière il se situe. Comparer le résultat à celui de
l'élève~$3$ et attribuer l'écart à l'une des deux variables.
\end{exercice}
